\sloppy

% Настройки стиля ГОСТ 7-32
% Для начала определяем, хотим мы или нет, чтобы рисунки и таблицы нумеровались в пределах раздела, или нам нужна сквозная нумерация.
% \EqInChapter % формулы будут нумероваться в пределах раздела
% \TableInChapter % таблицы будут нумероваться в пределах раздела
% \PicInChapter % рисунки будут нумероваться в пределах раздела

% Для закладок PDF: текст как есть (без \MakeUppercase) — нужно до hyperref
\makeatletter
\providecommand{\MakeUppercaseUnsupportedInPdfStrings}[1]{#1}
\makeatother

% Добавляем гипертекстовое оглавление в PDF
\usepackage[
      bookmarks=true, colorlinks=true, unicode=true,
      urlcolor=black,linkcolor=black, anchorcolor=black,
      citecolor=black, menucolor=black, filecolor=black,
]{hyperref}

\urlstyle{same}

% Автоматическая сортировка номеров источников по возрастанию и сворачивание
% подряд идущих в диапазон: [3,1,2] -> [1--3] (sort и compress включены
% в пакете cite по умолчанию)
\usepackage[nospace]{cite}
\renewcommand{\citedash}{\textendash}% длинное тире вместо двойного дефиса

% Патч пакета cite: по умолчанию пакет сжимает в диапазон только
% последовательности из 3+ номеров ([5,6,7] -> [5-7]), но пары
% оставляет как [5,6]. Переопределяем \@compress@cite, чтобы пары
% также сворачивались: [5,6] -> [5-6].
\makeatletter
\def\@compress@cite#1#2#3{%
  \ifnum\@cite@incr=\z@%
    \advance\@tempcnta\@cclvi%
    \ifnum #1=\@tempcnta%
       \expandafter\def\expandafter\@h@ld\expandafter{\citedash\@cite@out{#3}}%
       \mathchardef\@cite@incr=\@cclvi%
    \else%
       \advance\@tempcnta-\@cclv%
       \ifnum #1=\@tempcnta%
         \expandafter\def\expandafter\@h@ld\expandafter{\citedash\@cite@out{#3}}%
         \mathchardef\@cite@incr=\@ne%
       \else%
         \@citea \@cite@out{#3}%
       \fi%
    \fi%
  \else%
    \advance\@tempcnta\@cite@incr%
    \ifnum #1=\@tempcnta%
       \def\@h@ld{\citedash \@cite@out{#3}}%
    \else%
       \@h@ld \@citea \@cite@out{#3}%
       \let\@h@ld\@empty%
       \mathchardef\@cite@incr=\z@%
    \fi%
  \fi%
  \@tempcnta#1\let\@citea\citepunct%
}
\makeatother

\usepackage{pdfpages}

% Изменение начертания шрифта --- после чего выглядит таймсоподобно.
% apt-get install scalable-cyrfonts-tex

% \IfFileExists{cyrtimes.sty}
%     {
% \usepackage{cyrtimespatched}
%     }
%     {
%         % А если Times нету, то будет CM...
%     }

\usepackage{graphicx}   % Пакет для включения рисунков
\graphicspath{{images/}}

% С такими оно полями оно работает по-умолчанию:
% \RequirePackage[left=20mm,right=10mm,top=20mm,bottom=20mm,headsep=0pt]{geometry}
% Если вас тошнит от поля в 10мм --- увеличивайте до 20-ти, ну и про переплет не забывайте:

% Пакет Tikz
\usepackage{tikz}
\usetikzlibrary{arrows,positioning,shadows}

% Произвольная нумерация списков.
\usepackage{enumerate}

% ячейки в несколько строчек
\usepackage{multirow}

% itemize внутри tabular
\usepackage{paralist,array}

% Таблицы
\usepackage{array} % расширенные возможности для работы с таблицами
\usepackage{tabularx} % автоматический подбор ширины столбцов
\usepackage{dcolumn} % выравнивание чисел по разделителю

\usepackage{setspace} % полуторный интервал
\setstretch{1.5}

\newenvironment{compactlist}{%
      \renewcommand{\labelenumi}{\hspace{13mm}\arabic{enumi}.}
      \setlength{\topsep}{0pt}%
      \setlength{\partopsep}{0pt}%
      \setlength{\itemsep}{0pt}%
      \setlength{\parsep}{0pt}%
      \begin{enumerate}%
            }{%
      \end{enumerate}%
}

\makeatletter
\renewcommand*\l@section{\@dottedtocline{1}{0em}{2.3em}}
\renewcommand*\l@subsection{\@dottedtocline{2}{0em}{3.2em}}
\renewcommand*\l@subsubsection{\@dottedtocline{3}{0em}{4.1em}}
\makeatother
\setcounter{tocdepth}{2}
\setcounter{secnumdepth}{2}


\usepackage{caption}
\usepackage{ragged2e}

% Переносы слов в таблицах (колонки с \RaggedRight)
\newcolumntype{L}{>{\RaggedRight\arraybackslash}X}
\newcolumntype{C}{>{\Centering\arraybackslash}X}
\newcolumntype{P}[1]{>{\RaggedRight\arraybackslash}p{#1}}

\DeclareCaptionJustification{centerednohyphen}{\Centering\hyphenpenalty=10000\exhyphenpenalty=10000}
\DeclareCaptionLabelSeparator{shortdash}{ \textendash\ }

\usepackage{caption}

\usepackage{float}
\floatplacement{table}{H}
\setlength{\floatsep}{0pt}
\setlength{\textfloatsep}{0pt}
\setlength{\intextsep}{7pt}

\captionsetup{aboveskip=7pt,belowskip=0pt,justification=centerednohyphen,singlelinecheck=false,labelsep=shortdash}
\captionsetup[table]{justification=RaggedRight,singlelinecheck=false}
\captionsetup[figure]{justification=centerednohyphen,singlelinecheck=false}
\captionsetup[listing]{justification=RaggedRight,singlelinecheck=false, position=top}
\captionsetup[listing]{position=top}

\usepackage{chngcntr}

\usepackage{mdframed}
\usepackage{xcolor}

\mdfsetup{
      linecolor=black,
      linewidth=1pt,
      roundcorner=5pt,
      innerleftmargin=5pt,
      innerrightmargin=5pt,
      innertopmargin=5pt,
      innerbottommargin=5pt,
      skipabove=0pt,
      skipbelow=0pt,
      leftmargin=0pt,
      rightmargin=0pt
}

% tcolorbox: поддерживает breakable и скругленные углы одновременно (mdframed с roundcorner на TikZ не переносится через страницу).
\usepackage[breakable,skins]{tcolorbox}

\usepackage[chapter]{minted}

\captionsetup[listing]{
      position=top,
      justification=RaggedRight,
      singlelinecheck=false
}

\floatstyle{plaintop}
\restylefloat{listing}

\setminted{
      linenos=true,           % включить номера строк
      numberblanklines=false, % не нумеровать пустые строки
      numbersep=5pt,          % расстояние между номерами строк и кодом
      xleftmargin=2.5em,      % отступ слева, чтобы номера не вылезали за границы
      fontsize=\fontsize{12pt}{14pt}\selectfont, % ГОСТ: листинги — 12pt; моноширинный Courier New задан через \setmonofont в G2-105.sty
      breakanywhere=true,
      breakautoindent=true,
      breaklines=true,
      tabsize=4,              % заменяет символы табуляции (^^I) на 4 пробела
}

% Номера строк листингов — Courier 12pt (по умолчанию fancyvrb даёт \rmfamily\tiny)
\renewcommand{\theFancyVerbLine}{%
      \ttfamily\fontsize{12pt}{14pt}\selectfont\arabic{FancyVerbLine}}

\renewcommand{\listingscaption}{Листинг}
\renewcommand{\listoflistingscaption}{Листинги}

% codelisting: листинг с переносом через страницу. Использование: \begin{codelisting}{подпись}{label} \begin{minted}{lang}...\end{minted} \end{codelisting}
\newenvironment{codelisting}[2]{%
      \par\vspace{5pt}%
      \noindent\captionof{listing}{#1}\label{#2}%
      \nopagebreak[4]%
      \begin{tcolorbox}[%
                  breakable,%
                  enhanced,%
                  colback=white,%
                  colframe=black,%
                  arc=5pt,%
                  boxrule=0.5pt,%
                  left=5pt, right=5pt, top=5pt, bottom=5pt,%
            ]%
            }{%
      \end{tcolorbox}%
      \vspace{5pt}%
}


\makeatletter
\let\oldtableofcontents\tableofcontents
\renewcommand{\tableofcontents}{%
      \begingroup
      \hyphenpenalty=10000
      \exhyphenpenalty=10000
      \oldtableofcontents
      \endgroup
}
\makeatother

\makeatletter
\def\@hangfrom#1{\setbox\@tempboxa\hbox{{#1}}%
      \hangindent 0pt%\wd\@tempboxa
      \noindent\box\@tempboxa}
\makeatother

\usepackage{etoolbox}
\AtEndEnvironment{listing}{\vspace{-10pt}}
\AtEndEnvironment{table}{\vspace{-10pt}}
\AtEndEnvironment{center}{\vspace{-14pt}}
\AtEndEnvironment{figure}{\vspace{-14pt}}

% Внутри таблиц используется одинарный межстрочный интервал по ГОСТ 7.32-2017
\AtBeginEnvironment{tabular}{\setstretch{1}}
\AtBeginEnvironment{tabularx}{\setstretch{1}}

% Листинги кода — одинарный межстрочный интервал по ГОСТ 7.32-2017
\AtBeginEnvironment{minted}{\setstretch{1}}


\usepackage{fontspec}

% Указываем шрифты для основного текста и для кириллицы:
\setmainfont{Times New Roman} % если установлен в системе
\newfontfamily\cyrillicfont{Times New Roman} % для кириллического текста

% Убирает точку после номера главы
\makeatletter
\renewcommand{\thesection}{\thechapter.\arabic{section}}
\renewcommand{\thesubsection}{\thesection.\arabic{subsection}}
\renewcommand{\@seccntformat}[1]{%
      {\fontsize{14pt}{21pt}\selectfont\bfseries \csname the#1\endcsname}\hspace{1em}%
}
\makeatother

\newcommand{\appendixchapter}[3]{%
  \chapter*{{\fontsize{14pt}{21pt}\selectfont\MakeUppercase{ПРИЛОЖЕНИЕ #1}}}%
  \phantomsection
  \addcontentsline{toc}{chapter}{ПРИЛОЖЕНИЕ #1}%
  \label{#3}%
  \begin{center}
    {\fontsize{14pt}{21pt}\selectfont #2}
  \end{center}
}
